% Chapter 22: Compliance and Intention-to-Treat Analysis
% A First Course in Causal Inference - Peng Ding
% Draft V2 - Final Version (after ch22-writer-2 review)

\chapter{混合分布拆解与工具变量不等式}\label{ch:22}

\section{引言：从随机实验到现实世界}\label{sec:22-intro}

在前面几章中，我们假设随机化实验的治疗分配能够被完美执行。然而，现实世界中的实验往往面临一个棘手的问题：\textbf{依从性（compliance）}问题。

想象一项药物临床试验。研究人员将患者随机分配到治疗组或对照组，但有些患者可能忘记服药，有些可能自行停药，还有些对照组患者可能通过其他途径获得了该药物。这种"分配与实际治疗不符"的现象，就是依从性问题的核心。

正如计量经济学家 \citet{AngristImbensRubin:1996} 所指出的：

\textit{``The essential problem of causal inference is that we can observe at most one of the potential outcomes for any unit.''}

在依从性问题的语境下，这句话有了更深的含义：我们不仅无法同时观察同一个体的治疗组和对照组结果，还可能无法控制该个体实际接受的治疗。

\textbf{本章的核心问题}：当治疗分配与实际治疗不一致时，我们如何估计治疗效果？

\textbf{历史脉络}：依从性问题的系统研究始于 \citet{Zelen:1979} 的开创性工作，随后被 \citet{ImbensAngrist:1994} 和 \citet{AngristImbensRubin:1996} 发展到潜在结果框架下。本章遵循这一脉络，从实验设计视角介绍工具变量方法。

\section{鼓励设计与依从性问题}\label{sec:22-encouragement-design}

\subsection{基本设定与记号}

考虑一个包含 $n$ 个单元的完全随机实验。记号如下：

\begin{itemize}
  \item $Z_i \in \{0, 1\}$：治疗分配（1 = 治疗组，0 = 对照组）
  \item $D_i \in \{0, 1\}$：实际接受的治疗（1 = 接受治疗，0 = 未接受治疗）
  \item $Y_i$：关注的结果
\end{itemize}

当存在某些单元满足 $Z_i \neq D_i$ 时，\textbf{依从性问题（noncompliance problem）}就出现了。这种情况在涉及人类被试的实验中尤为常见——实验人员无法强迫被试接受治疗，只能"鼓励"他们。

\textbf{潜在结果记号}：每个单元有四个潜在结果
\[
\{D_i(1), D_i(0), Y_i(1), Y_i(0)\},
\]
其中 $D_i(z)$ 是在治疗分配 $z$ 下会接受的治疗，$Y_i(z)$ 是在治疗分配 $z$ 下的结果。观测值为：
\[
D_i = Z_i D_i(1) + (1 - Z_i) D_i(0), \quad Y_i = Z_i Y_i(1) + (1 - Z_i) Y_i(0).
\]

\subsection{随机化假设与可识别性}

我们首先假设完全随机化且忽略协变量。

\begin{Assumption}[随机化]\label{asm:22-randomization}
治疗分配 $Z$ 与潜在结果独立：
\[
Z \perp\!\!\!\perp \{D(1), D(0), Y(1), Y(0)\}.
\]
\end{Assumption}

随机化允许我们识别治疗分配对 $D$ 和 $Y$ 的平均因果效应：\footnote{推导见附录 \cref{sec:appendix-22-itt-identification}}
\[
\tau_D = \mathbb{E}\{D(1) - D(0)\} = \mathbb{E}(D \mid Z = 1) - \mathbb{E}(D \mid Z = 0),
\]
\[
\tau_Y = \mathbb{E}\{Y(1) - Y(0)\} = \mathbb{E}(Y \mid Z = 1) - \mathbb{E}(Y \mid Z = 0).
\]

我们可以用简单的差值均值估计量 $\hat\tau_D$ 和 $\hat\tau_Y$ 来估计 $\tau_D$ 和 $\tau_Y$。

\subsection{意向治疗分析（ITT）}

报告 $\hat\tau_Y$ 及其标准误的分析称为\textbf{意向治疗分析（intention-to-treat, ITT analysis）}。\footnote{ITT 分析的名称来源于临床试验，其核心思想是按照原始分配而非实际接受的治疗进行分析。}

ITT 估计的是\textbf{治疗分配}对结果的影响，而非\textbf{实际接受的治疗}对结果的影响。这是一个重要的区别：

\[
\underbrace{\hat\tau_Y}_{\text{ITT 估计量}} \xrightarrow{p} \underbrace{\tau_Y}_{\text{治疗分配的因果效应}},
\]
但我们真正关心的科学问题往往是：
\[
\underbrace{\tau_{\text{treatment}}}_{\text{实际治疗的因果效应}} \neq \underbrace{\tau_Y}_{\text{治疗分配的因果效应}}.
\]

这就是依从性问题带来的核心挑战。

\section{潜在依从性类型与依从性因果效应}\label{sec:22-latent-compliance}

\subsection{四类依从性状态}\label{sec:22-compliance-types}

遵循 \citet{ImbensAngrist:1994} 的开创性工作，我们根据潜在治疗接受值 $\{D_i(1), D_i(0)\}$ 对总体进行分层。由于 $D$ 是二值的，共有四种可能的组合：

\begin{Definition}[依从性类型]\label{def:22-compliance-types}
令 $U_i = \{D_i(1), D_i(0)\}$，则：
\[
U_i =
\begin{cases}
\text{a（always taker）}, & \text{若 } D_i(1) = 1 \text{ 且 } D_i(0) = 1; \\
\text{c（complier）}, & \text{若 } D_i(1) = 1 \text{ 且 } D_i(0) = 0; \\
\text{d（defier）}, & \text{若 } D_i(1) = 0 \text{ 且 } D_i(0) = 1; \\
\text{n（never taker）}, & \text{若 } D_i(1) = 0 \text{ 且 } D_i(0) = 0.
\end{cases}
\]
\end{Definition}

\textbf{直观理解}：
\begin{itemize}
  \item \textbf{Always takers}：无论分配如何，都会接受治疗
  \item \textbf{Never takers}：无论分配如何，都不会接受治疗
  \item \textbf{Compliers}：服从分配——分配到治疗组则接受，分配到对照组则不接受
  \item \textbf{Defiers}：反着来——分配到治疗组反而不接受，分配到对照组反而接受
\end{itemize}

由于我们无法同时观察 $D_i(1)$ 和 $D_i(0)$，类型 $U_i$ 是\textbf{潜变量}。

\subsection{平均因果效应的分解}\label{sec:22-decomposition}

基于潜在类型 $U$，利用全概率公式可以将 $\tau_Y$ 分解为四项：\footnote{完整推导见附录 \cref{sec:appendix-22-tau-y-decomposition}}
\[
\begin{aligned}
\tau_Y = & \; \mathbb{E}\{Y(1) - Y(0) \mid U = \text{a}\} \Pr(U = \text{a}) \\
& + \mathbb{E}\{Y(1) - Y(0) \mid U = \text{c}\} \Pr(U = \text{c}) \\
& + \mathbb{E}\{Y(1) - Y(0) \mid U = \text{d}\} \Pr(U = \text{d}) \\
& + \mathbb{E}\{Y(1) - Y(0) \mid U = \text{n}\} \Pr(U = \text{n}).
\label{eq:22-tau-y-decomposition}
\end{aligned}
\]

这说明 $\tau_Y$ 是四个潜在亚组因果效应的加权平均。

类似地，$\tau_D$ 的分解为：\footnote{推导见附录 \cref{sec:appendix-22-tau-d-decomposition}}
\[
\begin{aligned}
\tau_D = & \; \mathbb{E}\{D(1) - D(0) \mid U = \text{a}\} \Pr(U = \text{a}) \\
& + \mathbb{E}\{D(1) - D(0) \mid U = \text{c}\} \Pr(U = \text{c}) \\
& + \mathbb{E}\{D(1) - D(0) \mid U = \text{d}\} \Pr(U = \text{d}) \\
& + \mathbb{E}\{D(1) - D(0) \mid U = \text{n}\} \Pr(U = \text{n}) \\
= & \; 0 \cdot \Pr(U = \text{a}) + 1 \cdot \Pr(U = \text{c}) + (-1) \cdot \Pr(U = \text{d}) + 0 \cdot \Pr(U = \text{n}) \\
= & \; \Pr(U = \text{c}) - \Pr(U = \text{d}).
\label{eq:22-tau-d-decomposition}
\end{aligned}
\]

\section{非参数识别}\label{sec:22-identification}

上述分解表明，仅凭观测数据无法唯一识别 $\tau_Y$——我们有四个未知数（四个亚组的因果效应）但只有两个方程（$\tau_Y$ 和 $\tau_D$ 的识别）。为了进一步识别，我们需要额外的假设。

\subsection{单调性假设}

\begin{Assumption}[单调性]\label{asm:22-monotonicity}
$\Pr(U = \text{d}) = 0$，或等价地，$D_i(1) \geq D_i(0)$ 对所有单元成立。即，\textbf{不存在 defiers}。
\end{Assumption}

\textbf{何时成立？} 单调性假设在\textbf{单侧依从性问题（one-sided noncompliance）}中自动成立：当对照组被试无法获得治疗时，$D_i(0) = 0$ 对所有单元成立，从而 $D_i(1) \geq D_i(0)$ 必然成立。

\textbf{可检验的含义}：在随机化下，单调性假设有一个可检验的含义：\footnote{推导见附录 \cref{sec:appendix-22-monotonicity-implication}}
\[
\Pr(D = 1 \mid Z = 1) \geq \Pr(D = 1 \mid Z = 0).
\label{eq:22-monotonicity-inequality}
\]

注意：\cref{asm:22-monotonicity} 比 \cref{eq:22-monotonicity-inequality} 更强。前者限制每个个体的 $D_i(1)$ 和 $D_i(0)$，后者只限制平均值。当可检验的不等式成立时，我们无法用数据反驳单调性假设。

在单调性假设下，\cref{eq:22-tau-d-decomposition} 简化为：
\[
\tau_D = \Pr(U = \text{c}) \equiv \pi_c.
\label{eq:22-tau-d-monotonicity}
\]

这意味着 compliers 的比例 $\pi_c$ 等于治疗分配对实际治疗的平均因果效应——这是一个可以识别的数量。

\subsection{排斥限制假设}

\begin{Assumption}[排斥限制]\label{asm:22-exclusion-restriction}
对于 always takers（$U = \text{a}$）和 never takers（$U = \text{n}$），有 $Y_i(1) = Y_i(0)$。
\end{Assumption}

\textbf{直观理解}：排斥限制要求治疗分配\textbf{只能通过实际接受的治疗}影响结果，不能有其他直接通路。在双盲随机对照试验中，这是生物学上合理的——如果治疗分配没有改变实际接受的治疗，它也不会改变结果。

\textbf{何时可能被违反？} 当治疗分配有直接影响结果的通路时（而非通过实际接受的治疗）。例如，在非双盲试验中，被试可能因为知道自己被分配到治疗组而改变行为。

在排斥限制假设下，\cref{eq:22-tau-y-decomposition} 中的第一项和最后一项变为零：
\[
\tau_Y = \mathbb{E}\{Y(1) - Y(0) \mid U = \text{c}\} \Pr(U = \text{c}).
\label{eq:22-tau-y-under-exclusion}
\]

结合 \cref{eq:22-tau-d-monotonicity} 和 \cref{eq:22-tau-y-under-exclusion}，我们得到识别结果。

\subsection{主要识别定理}

\begin{Theorem}[CACE/LATE 的识别]\label{th:22-cace-identification}
在 \cref{asm:22-randomization}、\cref{asm:22-monotonicity} 和 \cref{asm:22-exclusion-restriction} 下，如果 $\tau_D \neq 0$，则
\[
\mathbb{E}\{Y(1) - Y(0) \mid U = \text{c}\} = \frac{\tau_Y}{\tau_D}.
\label{eq:22-cace-formula}
\]
\end{Theorem}

\textbf{证明思路}：由 \cref{eq:22-tau-y-under-exclusion} 和 \cref{eq:22-tau-d-monotonicity} 可直接得到。\footnote{完整证明见附录 \cref{sec:appendix-22-cace-proof}}。

\cref{th:22-cace-identification} 的含义是：在上述三个假设下，$\tau_Y / \tau_D$ 识别了 compliers 的平均因果效应。

\textbf{一个哲学性的观察}：识别结果只涉及 compliers——那些"遵医嘱"的人。这意味着我们无法直接得知治疗对 always takers、never takers 或 defiers 的效果。这既是局限，也是优势——它允许我们专注于一个明确定义的亚组。

\section{依从性平均因果效应（CACE/LATE）}\label{sec:22-cace}

\subsection{定义}\label{sec:22-cace-definition}

\begin{Definition}[CACE 或 LATE]\label{def:22-cace}
\textbf{依从性平均因果效应（complier average causal effect, CACE）}或\textbf{局部平均治疗效应（local average treatment effect, LATE）}定义为：
\[
\tau_c = \mathbb{E}\{Y(1) - Y(0) \mid U = \text{c}\}.
\]
\end{Definition}

CACE 有两种等价形式：
\[
\tau_c = \mathbb{E}\{Y(1) - Y(0) \mid D(1) = 1, D(0) = 0\} = \mathbb{E}\{Y(1) - Y(0) \mid D(1) > D(0)\}.
\]

基于 \cref{def:22-cace}，\cref{th:22-cace-identification} 可以写成：
\[
\tau_c = \frac{\tau_Y}{\tau_D}.
\label{eq:22-cace-ratio}
\]

即，\textbf{CACE/LATE 等于结果平均因果效应与治疗接受平均因果效应之比}。

在完全随机化下，进一步有：
\[
\tau_c = \frac{\mathbb{E}(Y \mid Z = 1) - \mathbb{E}(Y \mid Z = 0)}{\mathbb{E}(D \mid Z = 1) - \mathbb{E}(D \mid Z = 0)}.
\label{eq:22-cace-observed}
\]

因此，在完全随机化、单调性和排斥限制下，我们可以\textbf{非参数地识别 CACE}——它等于观测数据中结果差值均值与治疗接受差值均值之比。

\section{估计}\label{sec:22-estimation}

\subsection{Wald 估计量}\label{sec:22-wald-estimator}

基于 \cref{eq:22-cace-observed}，我们可以用简单的比值估计 CACE：
\[
\hat{\tau}_c = \frac{\hat{\tau}_Y}{\hat{\tau}_D},
\label{eq:22-wald-estimator}
\]
其中
\[
\hat{\tau}_Y = \frac{1}{n_1} \sum_{i: Z_i=1} Y_i - \frac{1}{n_0} \sum_{i: Z_i=0} Y_i, \quad
\hat{\tau}_D = \frac{1}{n_1} \sum_{i: Z_i=1} D_i - \frac{1}{n_0} \sum_{i: Z_i=0} D_i.
\]

这就是著名的 \textbf{Wald 估计量}（\citealp{Wald:1940}），也称为 \textbf{工具变量估计量}。在这个语境下，$Z$ 作为 $D$ 的工具变量。

\textbf{直观理解}：Wald 估计量可以理解为"缩放的" ITT 估计量。$\hat{\tau}_Y$ 估计治疗分配的效果，$\hat{\tau}_D$ 估计治疗分配对实际治疗的"转化率"（即 compliers 的比例）。二者的比值估计了对于那些"被转化"的 compliers，治疗的实际效果。

\subsection{方差估计}\label{sec:22-variance-estimation}

为了进行统计推断，我们需要估计 $\hat{\tau}_c$ 的方差。\footnote{渐近方差公式的推导见附录 \cref{sec:appendix-22-asymptotic-variance}}。

基本思想是将 $\hat{\tau}_c$ 在真实值附近线性化：\footnote{详细推导见附录 \cref{sec:appendix-22-wald-delta-method}}
\[
\hat{\tau}_c - \tau_c \approx \frac{\hat{\tau}_Y - \tau_c \hat{\tau}_D}{\tau_D}.
\]

定义调整后的结果 $A_i = Y_i - \tau_c D_i$，则上式近似为 $\hat{\tau}_A / \tau_D$，其中 $\hat{\tau}_A$ 是调整后结果的差值均值。

因此，$\hat{\tau}_c$ 的渐近方差约为 $\hat{V}_A / \tau_D^2$，其中 $\hat{V}_A$ 是调整后结果的 Neyman 型方差估计量：

\[
\hat{V}_A = \frac{\hat{S}_A^2(1)}{n_1} + \frac{\hat{S}_A^2(0)}{n_0},
\label{eq:22-adjusted-variance}
\]

其中 $\hat{S}_A^2(1)$ 和 $\hat{S}_A^2(0)$ 分别是治疗组和对照组中调整后结果的样本方差。

最终，我们使用
\[
\hat{V}_{\hat{\tau}_c} = \frac{\hat{V}_A}{\hat{\tau}_D^2}
\label{eq:22-cace-variance}
\]
作为 $\hat{\tau}_c$ 的方差估计量。

基于此，我们可以构建 $\tau_c$ 的 $1-\alpha$ 置信区间：
\[
\hat{\tau}_c \pm z_{1-\alpha/2} \sqrt{\hat{V}_{\hat{\tau}_c}}.
\label{eq:22-cace-ci}
\]

\section{弱工具变量}\label{sec:22-weak-iv}

在实际研究中，我们常常面临一个尴尬的局面：工具变量的"强度"并不总能得到保证。想象一下，即使在真实的随机实验中，由于各种现实因素，只有极少数被试（比如说 5\%）改变了他们的行为。在这种情况下，我们能否信赖 Wald 估计量？

\subsection{问题诊断}\label{sec:22-weak-iv-problem}

即使 $\tau_D > 0$（即 compliers 确实存在），$\hat{\tau}_D$ 仍有可能接近甚至等于零。这会导致 Wald 估计量的方差急剧增大——实际上，当 $\tau_D \to 0$ 时，$\text{var}(\hat{\tau}_c) \to \infty$。\footnote{这与经典统计学中的 Fieller-Creasy 问题类似。}

这就是\textbf{弱工具变量（weak instrumental variable）}问题。\citet{StaigerStock:1997} 的经典研究表明，当工具变量弱时，传统的统计推断方法会失效。在这种情况下：
\begin{itemize}
  \item $\hat{\tau}_c$ 存在有限样本偏
  \item 渐近分布不再是正态
  \item 基于正态近似的置信区间覆盖性质很差
\end{itemize}

\textbf{诊断指标}：一个常用的诊断是 $F$ 统计量或 $\hat{\tau}_D$ 的绝对值。经验法则：当 $\hat{\tau}_D$ 很小时，应谨慎对待 Wald 估计量。

\subsection{FAR 置信区间}\label{sec:22-far-ci}

从估计角度，当工具变量弱时，我们建议关注\textbf{置信区间}而非点估计。从检验角度，有一个简单的解决方案——由于 $\tau_c = \tau_Y / \tau_D$，原假设 $\tau_c = 0$ 等价于 $\tau_Y = 0$。

更精细的方法是构建\textbf{Fieller-Anderson-Rubin（FAR）置信区间}。其思想是通过反转一系列假设检验来构建 $\tau_c$ 的置信域。\footnote{FAR 置信区间的详细讨论见附录 \cref{sec:appendix-22-far-ci}。}

对于给定的 $b$，考虑原假设 $H_0(b): \tau_c = b$。这等价于 $\tau_Y - b\tau_D = 0$，即调整后的结果 $A(b)_i = Y_i - b D_i$ 的平均因果效应为零。

构造 Wald 型检验并反转，可以得到置信域：
\[
\left\{ b: \frac{\hat{\tau}_A^2(b)}{\hat{V}_A(b)} \leq z_{1-\alpha/2}^2 \right\},
\label{eq:22-far-set}
\]
其中 $\hat{\tau}_A(b)$ 和 $\hat{V}_A(b)$ 是调整后结果 $A(b)_i = Y_i - b D_i$ 的点估计和方差估计。

FAR 置信域可能是单个区间、两个不相连的区间、空集或整个实数线。当工具变量强时，它收敛到标准 Wald 置信区间；当工具变量弱时，它提供更可靠的推断。

\section{协变量调整}\label{sec:22-covariates}

到目前为止我们都假设没有协变量。加入协变量可以提高估计的效率和稳健性。

\subsection{完全随机化下的协变量调整}

假设条件随机化成立：
\[
Z \perp \{D(1), D(0), Y(1), Y(0), X\}.
\]

我们可以使用 \citet{Lin:2013} 的协变量调整估计量。对于结果和和治疗接受，分别拟合包含协变量的回归模型：
\[
D_i = \beta_0 + \beta_1 Z_i + \beta_2 X_i + \beta_3 Z_i X_i + \epsilon_i,
\]
\[
Y_i = \alpha_0 + \alpha_1 Z_i + \alpha_2 X_i + \alpha_3 Z_i X_i + \eta_i,
\]
取 $\hat{\tau}_{D,\text{L}} = \hat{\beta}_1$ 和 $\hat{\tau}_{Y,\text{L}} = \hat{\alpha}_1$，则调整后的 CACE 估计量为：
\[
\hat{\tau}_{c,\text{L}} = \frac{\hat{\tau}_{Y,\text{L}}}{\hat{\tau}_{D,\text{L}}}.
\label{eq:22-lin-cace}
\]

渐近方差可以通过 Bootstrap 近似。

\section{实例：JOBS II 培训项目}\label{sec:22-application}

\textbf{研究背景}：JOBS II 是一个随机化现场实验，研究职业培训对失业工人求职自我效能的影响。\citet{JuddKenny:1981} 原始设计比较复杂，这里使用 \citet{ImbensKlein:2015} 的简化版本。

\textbf{数据}：来自 \texttt{jobsdata.csv}。关键变量：
\begin{itemize}
  \item \texttt{treat}：是否被随机分配参加培训项目（$Z$）
  \item \texttt{comply}：是否实际参加了培训项目（$D$）
  \item \texttt{jobseek}：求职自我效能得分（$Y$，1-5 分制）
\end{itemize}

\textbf{描述性统计}：\footnote{详细数据分析代码见附录 \cref{sec:appendix-22-jobs-analysis}}

\begin{Example}[JOBS II 数据的 IV 分析]\label{ex:22-jobs-iv}
我们使用 Wald 估计量分析 JOBS II 数据。结果如下：

\begin{center}
\begin{tabular}{lcccc}
\hline
方法 & 估计量 & 标准误 & 95\% CI 下限 & 95\% CI 上限 \\
\hline
Delta 方法 & 0.109 & 0.081 & -0.050 & 0.268 \\
Bootstrap & 0.109 & 0.083 & -0.054 & 0.271 \\
含协变量调整 & 0.118 & 0.082 & -0.042 & 0.278 \\
\hline
\end{tabular}
\end{center}

\textbf{解读}：点估计表明，对于 compliers，参加培训项目使求职自我效能得分提高约 0.11 分（满分 5 分）。但 95\% 置信区间包含零，这意味着在常规显著性水平下，我们无法拒绝"培训对 compliers 无效果"的原假设。
\end{Example}

\section{从识别到解释：CACE 的深层含义}\label{sec:22-cace-interpretation}

到目前为止，我们建立了 CACE 的识别和估计框架。但还有一个关键问题：CACE 到底代表什么？

\cref{def:22-cace} 中定义的 $\tau_c$ 有两种等价解释：

\begin{enumerate}
  \item \textbf{治疗分配的因果效应（among compliers）}：$\tau_c = \mathbb{E}\{Y(1) - Y(0) \mid U = \text{c}\}$ 是治疗分配对 compliers 结果的平均因果效应。
  \item \textbf{实际治疗的因果效应（among compliers）}：对于 compliers，有 $D(1) = 1$ 和 $D(0) = 0$，因此 $D = Z$。所以 $\tau_c$ 也是实际接受的治疗对 compliers 结果的因果效应。
\end{enumerate}

第二种解释更具科学吸引力——它回答了"对于那些会遵医嘱的人，治疗真的有效果吗？"这个问题。

\textbf{局限性}：CACE 只告诉我们关于 compliers 的信息，我们无法直接得知治疗对 always takers、never takers 或 defiers 的效果。如果这些亚组与 compliers 有本质不同的治疗效果，CACE 可能无法代表治疗对整个人群的效果。

\section{本章小结}\label{sec:22-summary}

本章介绍了依从性问题的框架，主要内容如下：

\begin{enumerate}
  \item \textbf{依从性问题的普遍性}：在涉及人类被试的实验中，治疗分配与实际治疗往往不一致。
  \item \textbf{四类依从性状态}：always takers、compliers、defiers、never takers。
  \item \textbf{关键假设}：随机化、单调性（无 defiers）、排斥限制（治疗分配只能通过实际治疗影响结果）。
  \item \textbf{识别}（\cref{th:22-cace-identification}）：在上述假设下，CACE = $\tau_Y / \tau_D$。
  \item \textbf{估计}：Wald 估计量 $\hat{\tau}_c = \hat{\tau}_Y / \hat{\tau}_D$。
  \item \textbf{弱工具变量问题}：当 compliers 比例很小时，Wald 估计量表现不佳，需要使用 FAR 置信区间。
  \item \textbf{协变量调整}：可以提高效率和稳健性。
\end{enumerate}

\section{习题}\label{sec:22-homework}

\begin{HomeworkProblem}[More detailed data analysis]\label{exr:22-1}

Examples \ref{eg::1-iv-inequalities} and \ref{eg::2-iv-inequalities} ignored the uncertainty in the estimates. Calculate the confidence intervals for the true parameters.




\end{HomeworkProblem}

\begin{HomeworkProblem}[Risk ratio for compliers]\label{exr:22-2}


With a binary outcome, we can define the risk ratio for compliers as
\[
\textsc{rr}_\cp = \frac{  \pr\{  Y(1) = 1\mid U = \cp \}  }{ \pr\{  Y(0) = 1\mid U = \cp \}   } .
\]

Show that under Assumptions \ref{assume::random}--\ref{assume::ER}, we can identify it by
\[
\textsc{rr}_\cp =  \frac{  E(DY\mid Z=1)  - E(DY\mid Z=0)  }{  E\{  (D-1)Y\mid Z=1 \}   - E\{ (D-1)Y\mid Z=0 \} }.
\]



Remark: Using Theorem \ref{thm::cace-mixture-components}, we can identify any comparisons between $E\{  Y(1) \mid U = \cp \} $ and $E\{  Y(0) \mid U = \cp \} $.


\end{HomeworkProblem}

\begin{HomeworkProblem}[Disentangle the mixtures: distributional results]\label{exr:22-3}


This problem extends Theorem \ref{thm::cace-mixture-components}. Define
\[
f_{zu}(y) = \pr\{  Y(z) = y  \mid U=u \},  \quad  (z=0,1;u=\at, \nt, \cp)
\]
as the density of $Y(z)$ for latent stratum $U=u$, and define
\[
g_{zd}(y) = \pr(Y = y\mid Z=z, D=d)
\]
as the density of the outcome within the observed group $(Z=z, D=d)$.
Exclusion restriction implies that $f_{1\nt}(y) = f_{0\nt} (y) $ and $f_{1\at}(y) = f_{0\at}(y) $. Let $f_{\nt}(y)$ and $ f_{\at} (y) $ denote them, respectively.



Prove  Theorem \ref{thm::cace-mixture-components-distribution} below.



\begin{Theorem}
\label{thm::cace-mixture-components-distribution}
Under Assumptions \ref{assume::random}--\ref{assume::ER}, we can identify the type-specific densities of the potential outcomes by
\begin{eqnarray*}
  f_{\nt}(y) &=& g_{10}(y) ,\\
  f_{\at} (y) &=& g_{01}(y),\\
f_{1\cp}(y) &=& \pi_\cp^{-1} \{   \pr(D=1\mid Z=1 )g_{11}(y)  -  \pr(D=1\mid Z=0) g_{01}(y)   \} , \\
f_{0\cp} (y) &=& \pi_\cp^{-1} \{  \pr(D=0\mid Z=0)  g_{00}(y)  -  \pr(D=0\mid Z=1)  g_{10}(y)    \} .
\end{eqnarray*}
\end{Theorem}









\end{HomeworkProblem}

\begin{HomeworkProblem}[Alternative form of Theorem \ref{thm::iv-inequality}]\label{exr:22-4}

The inequalities in \eqref{eq::iv-inequalities} can be re-written as
\begin{eqnarray*}
\pr(D=1, Y=y\mid Z=1) &\geq & \pr(D=1, Y=y\mid Z=0), \\
\pr(D=0, Y=y\mid Z=0) &\geq & \pr(D=0, Y=y\mid Z=1)
\end{eqnarray*}
for both $y=0,1$.



\end{HomeworkProblem}

\begin{HomeworkProblem}[IV inequalities for a general outcome]\label{exr:22-5}

For a general outcome $Y$, show that Assumptions \ref{assume::random}--\ref{assume::ER} imply
\begin{eqnarray*}
\pr(D=1, Y\geq y\mid Z=1) &\geq & \pr(D=1, Y\geq y\mid Z=0), \\
\pr(D=1, Y <  y\mid Z=1) &\geq & \pr(D=1, Y <  y\mid Z=0), \\
\pr(D=0, Y\geq y\mid Z=0) &\geq & \pr(D=0, Y\geq y\mid Z=1),\\
\pr(D=0, Y <  y\mid Z=0) &\geq & \pr(D=0, Y <  y\mid Z=1)
\end{eqnarray*}
for all $y$.

Remark:  \citet{imbens1997estimating} and \citet{kitagawa2015test} discussed similar results.  We can test the first inequality based on an analog of the Kolmogorov--Smirnov statistic:
\[
\text{KS}_1 = \max_{y} \Big |
\frac{  \sumn Z_i D_i  1(Y_i \leq y)  }{ \sumn Z_i D_i  } - \frac{  \sumn (1-Z_i) D_i  1(Y_i \leq y)  }{ \sumn (1-Z_i) D_i  }
\Big |.
\]



\end{HomeworkProblem}

\begin{HomeworkProblem}[Example for the IV inequalities]\label{exr:22-6}

Give an example in which all the IV inequalities hold and another example in which not all the IV inequalities hold. You need to specify the joint distribution of $(Z,D,Y)$ with binary variables.



\end{HomeworkProblem}

\begin{HomeworkProblem}[Violations of the key assumptions]\label{exr:22-7}

Theorem \ref{thm::CACE-identification} relies on randomization, monotonicity, and exclusion restriction. The latter two are not testable even in randomized experiments. When they are violated, the IV estimator no longer identifies the CACE. This problem gives two cases below, which are restatements of  Propositions 2 and 3 in \citet{angrist1996identification}. Recall $\pi_u = \pr(U=u)$ and $\tau_u = E\{  Y(1) - Y(0) \mid U=u \}$ for $u=\at, \nt, \cp, \df$.


\begin{Theorem}
\label{theorem::wald-violation-assumptions}
(a)
Under Assumptions  \ref{assume::random} and \ref{assume::Mono} without the exclusion restriction, we have
\[
\frac{ E(Y\mid Z=1) - E(Y\mid Z=0) }{ E(D\mid Z=1) - E(D \mid Z=0) } - \tau_\cp
= \frac{  \pi_\at \tau_\at + \pi_\nt \tau_\nt  }{  \pi_\cp } .
\]


(b)
Under Assumptions  \ref{assume::random} and \ref{assume::ER} without the monotonicity, we have
\[
\frac{ E(Y\mid Z=1) - E(Y\mid Z=0) }{ E(D\mid Z=1) - E(D \mid Z=0) } - \tau_\cp
= \frac{  \pi_\df (\tau_\cp + \tau_\df )  }{   \pi_\cp - \pi_\df   }.
\]
\end{Theorem}

Prove Theorem \ref{theorem::wald-violation-assumptions}.


\end{HomeworkProblem}

\begin{HomeworkProblem}[Problems of other analyses]\label{exr:22-8}

In the process of deriving the IV inequalities in Section \ref{sec::disentangle-mixtures}, we disentangled the mixture distributions by identifying the proportions of the latent strata as well as the conditional means of their potential outcomes. These results help to understand the drawbacks of other seemingly reasonable analyses. I review three estimators below and suppose Assumptions \ref{assume::random}--\ref{assume::ER} hold.


\begin{enumerate}
\item
The {\it as-treated analysis} compares the means of the outcomes among units receiving the treatment and control, yielding
\[
\tau_\textsc{at} = E(Y\mid D=1) - E(Y\mid D=0).
\]
Show that
\[
\tau_\textsc{at} = \frac{  \pi_\at \mu_\at + \pr(Z=1)\pi_\cp \mu_{1\cp}   }{   \pr(D=1) }
- \frac{  \pi_\nt \mu_\nt + \pr(Z=0) \pi_\cp \mu_{0\cp}   }{   \pr(D=0) } .
\]


\item
The {\it per-protocol analysis} compares the units that comply with the treatment assigned in treatment and control groups, yielding
\[
\tau_\textsc{pp} = E(Y\mid Z=1, D=1) - E(Y\mid Z=0, D=0).
\]
Show that
\[
\tau_\textsc{pp} =\frac{  \pi_\at \mu_\at + \pi_\cp \mu_{1\cp} }{  \pi_\at   + \pi_\cp }
- \frac{  \pi_\nt \mu_\nt + \pi_\cp \mu_{0\cp} }{  \pi_\nt   + \pi_\cp }.
\]


\item
We may also want to compare the outcomes among units receiving the treatment and control, conditioning on their treatment assignment, yielding
\begin{eqnarray*}
\tau_{Z=1} &=& E(Y\mid Z=1, D=1) - E(Y\mid Z=1, D=0),\\
\tau_{Z=0} &=& E(Y\mid Z=0, D=1) - E(Y\mid Z=0, D=0).
\end{eqnarray*}
Show that they reduce to
\[
\tau_{Z=1}  = \frac{  \pi_\at \mu_\at + \pi_\cp \mu_{1\cp} }{  \pi_\at   + \pi_\cp } - \mu_\nt,\quad
\tau_{Z=0} = \mu_\at - \frac{  \pi_\nt \mu_\nt + \pi_\cp \mu_{0\cp} }{  \pi_\nt   + \pi_\cp }.
\]
\end{enumerate}


\end{HomeworkProblem}

\begin{HomeworkProblem}[Bounds on the average causal effect on the whole population]\label{exr:22-9}


Extend the discussion in Section \ref{sec::disentangle-mixtures} based on the notation in Section \ref{eq::interpret-iv-double}. With the potential outcome $Y(d)$, define the average causal effect of the treatment received on the outcome as
\[
\delta = E\{  Y(d=1) - Y(d=0) \},
\]
and modify the definition of $\mu_{zu} $ as
\[
m_{du} = E\{  Y(d) \mid U = u\},\quad  (z=0,1;u=\at, \nt, \cp)
\]
due to the change of the notation.
They satisfy
\[
\delta = \sum_{u = \at, \nt, \cp}  \pi_u  (m_{1u} - m_{0u}).
\]

Section \ref{sec::disentangle-mixtures} identifies $\pi_\at$, $\pi_\nt$, $\pi_\cp$, $m_{1\at} = \mu_{1 \at}$, $m_{0\nt} = \mu_{0\nt }$, $m_{1\cp} = \mu_{1\cp}$ and $m_{0\cp} = \mu_{0\cp}$. But the data do not contain any information about $m_{0 \at}$ and $m_{1\nt }$. Therefore, we cannot identify $\delta$. With a bounded outcome, we can bound $\delta$.

\begin{Theorem}
\label{thm::bound-on-ACE}
Under Assumptions \ref{assume::Mono}--\ref{assume::ER-double-index} with a bounded outcome in $[\underline{y}, \overline{y}]$, we have $\underline{\delta} \leq  \delta  \leq \overline{\delta}$, where
\[
\underline{\delta} =  \delta'  - \overline{y} \pr(D=1\mid Z=0) + \underline{y} \pr(D=0\mid Z=1)
\]
and
\[
\overline{\delta} = \delta'  -   \underline{y} \pr(D=1\mid Z=0) +  \overline{y}  \pr(D=0\mid Z=1)
\]
with $\delta' = E(DY\mid Z=1) - E(Y-DY\mid Z=0) .$
\end{Theorem}

Prove Theorem  \ref{thm::bound-on-ACE}.


Remark: In the special case with a binary outcome, the bounds simplify to
 \[
\underline{\delta} =   E(DY\mid Z=1) - E( D+Y-DY\mid Z=0)
\]
and
\[
\overline{\delta} =  E(DY+1-D\mid Z=1) - E(Y-DY\mid Z=0) .
 \]


\end{HomeworkProblem}

\begin{HomeworkProblem}[One-sided noncompliance and statistical inference]\label{exr:22-10}

Consider a randomized encouragement design where the units assigned to the control have no access to the treatment. For unit $i$, let $Z_i$ be the binary treatment assigned, $D_i$ be the binary treatment received, and $Y_i$ be the outcome of interest. One-sided noncompliance happens when
\[
Z_i=0 \Longrightarrow  D_i = 0 \quad (i=1,\ldots, n).
\]
Suppose that Assumption \ref{assume::random} holds.


\begin{enumerate}
[(1)]
\item
Does monotonicity Assumption \ref{assume::Mono} hold in this case?
How many latent strata defined by $\{ D_i(1), D_i(0)  \}$ are there in this problem? How do we identify their proportions by the observed data distribution?

\item
State the assumption of exclusion restriction. Under exclusion restriction, show that $E\{  Y(z) \mid U = u \}$ can be identified by the observed data distributions.
Give the formulas for all possible values of $z$ and $u$. How do we identify the CACE in this case?



\item
If we observe pretreatment covariates $X_i$ for all units $i$, how do we use the covariate information to improve the estimation efficiency of the CACE?


 \item
 Under Assumption \ref{assume::random}, the exclusion restriction Assumption \ref{assume::ER} has testable implications, which are the IV inequalities for one-sided noncompliance. State the IV inequalities.

\item
\citet{sommer1991estimating} provided the following dataset:

 \begin{center}
\begin{tabular}{cccccc}
\hline
      &   \multicolumn{2}{c}{$Z=1$}  &  &  \multicolumn{2}{c}{$Z=0$}  \\  \cline{2-3} \cline{5-6}
      & $D=1$  &  $D=0$           &       &  $D=1$  &  $D=0$\\
 $Y=1$ &   9663    &     2385         &      &     0       &    11514       \\
 $Y=0$ &   12   &     34            &      &     0         &    74           \\
 \hline
\end{tabular}
 \end{center}

The treatment assigned $Z$ is whether or not the child was assigned to the vitamin A supplement, the treatment received indicator $D$ is whether or not the child received the vitamin A supplement, and the binary outcome $Y$ is the survival indicator. The original RCT was conducted in Indonesia.
Re-analyze the data.
 \end{enumerate}


Remark: \citet{bloom1984accounting} first discussed one-sided noncompliance and proposed the estimator $\hat\tau_\cp = \hat\tau_Y / \hat\tau_D$, which is sometimes called the Bloom estimator.  \citet{bloom1984accounting}'s notation is different from this chapter.





\end{HomeworkProblem}

\begin{HomeworkProblem}[One-sided noncompliance with partial adherence]\label{exr:22-11}

\citet[][Table 3]{sanders2021incorporating} reported the following data from an RCT aiming to estimate the efficacy of smoking cessation interventions among individuals with psychotic disorders.

\begin{center}
\begin{tabular}{cccc}
\hline
group assigned& treatment received &group size& \# positive outcomes \\
\hline
Control & None& 151 &25 \\
Treatment& None& 35 & 7 \\
Treatment& Partial& 42& 17 \\
Treatment& Full& 70& 40\\
\hline
\end{tabular}
\end{center}

Three tiers of treatment received are defined as follows: ``full'' treatment corresponds to attending all $8$ treatment sessions, ``partial'' corresponds to attending $5$ to $7$ sessions, and ``none'' corresponds to $<5$ sessions. The outcome is
defined as the binary indicator of smoking reduction of $50\%$ or greater relative to baseline, measured at three months.


In this problem, the treatment assignment $Z$ is binary but the treatment received $D$ takes three values $0, 0.5, 1$ for ``none'', ``partial'', and ``full.''  The three-leveled $D$ causes complications, but it can only be 0 under the control assignment. How many latent strata $U = \{ D(1), D(0) \} $ do we have in this problem? Can we identify their proportions?

How do we extend the exclusion restriction to this problem? What can be the causal effects of interest? Can we identify them?

Analyze the data based on the questions above.





\end{HomeworkProblem}

\begin{HomeworkProblem}[Recommended reading]\label{exr:22-12}

\citet{balke1997bounds} derived more general IV inequalities.


\end{HomeworkProblem}
